\documentclass[../../main]{subfiles}

\begin{document}

\section{Produto Cartesiano}
\label{sec:matematica:teoria-conjuntos:produto-cartesiano}

\begin{defn}[Par Ordenado]
	\label{def:matematica:teoria-conjuntos:produto-cartesiano:par-ordenado}

	Um par ordenado é uma coleção de dois objetos na qual a ordem dos
	componentes é relevante.

	Um par ordenado cujos componentes são (a) e (b) é denotado por

		[
			(a,b).
		]

\end{defn}

\begin{prop}[Igualdade de Pares Ordenados]
	\label{prop:matematica:teoria-conjuntos:produto-cartesiano:igualdade-pares}

	Para quaisquer objetos (a), (b), (c) e (d),

	[
			(a,b)=(c,d)
		]

	se, e somente se,

	[
			a=c
			\quad\text{e}\quad
			b=d.
		]

\end{prop}

\begin{proof}

	A igualdade de pares ordenados ocorre exatamente quando seus
	componentes correspondentes são iguais.

\end{proof}

\begin{defn}[Primeira Projeção]
	\label{def:matematica:teoria-conjuntos:produto-cartesiano:primeira-projecao}

	Seja

		[
			(a,b)
		]

	um par ordenado.

	Define-se sua primeira projeção por

		[
			\pi_1(a,b)=a.
		]

\end{defn}

\begin{defn}[Segunda Projeção]
	\label{def:matematica:teoria-conjuntos:produto-cartesiano:segunda-projecao}

	Seja

		[
			(a,b).
		]

	Define-se sua segunda projeção por

		[
			\pi_2(a,b)=b.
		]

\end{defn}

\begin{defn}[Produto Cartesiano]
	\label{def:matematica:teoria-conjuntos:produto-cartesiano:produto-cartesiano}

	Sejam (A) e (B) conjuntos.

	O produto cartesiano de (A) por (B), denotado por

		[
			A\times B,
		]

	é definido por

	[
	A\times B
	=
	{
	(a,b)
	\mid
	a\in A
	\land
	b\in B
	}.
	]

\end{defn}

\begin{ex}

	Se

	[
	A={1,2}
	\quad\text{e}\quad
	B={a,b},
	]

	então

	[
	A\times B
	=
	{
	(1,a),
	(1,b),
	(2,a),
	(2,b)
	}.
	]

\end{ex}

\begin{prop}[Caracterização do Produto Cartesiano]
	\label{prop:matematica:teoria-conjuntos:produto-cartesiano:caracterizacao}

	Para quaisquer conjuntos (A) e (B),

	[
			(x,y)\in A\times B
		]

	se, e somente se,

	[
			x\in A
			\quad\text{e}\quad
			y\in B.
		]

\end{prop}

\begin{proof}

	A equivalência segue diretamente da definição.

\end{proof}

\begin{prop}[Produto Cartesiano com o Conjunto Vazio]
	\label{prop:matematica:teoria-conjuntos:produto-cartesiano:vazio}

	Para qualquer conjunto (A),

	[
			A\times\varnothing
			=
			\varnothing
		]

	e

		[
			\varnothing\times A
			=
			\varnothing.
		]

\end{prop}

\begin{proof}

	Não existem pares ordenados cujos componentes pertençam
	simultaneamente aos conjuntos considerados.

\end{proof}

\begin{prop}[Monotonicidade do Produto Cartesiano]
	\label{prop:matematica:teoria-conjuntos:produto-cartesiano:monotonicidade}

	Se

		[
			A\subseteq C
			\quad\text{e}\quad
			B\subseteq D,
		]

	então

	[
	A\times B
	\subseteq
	C\times D.
	]

\end{prop}

\begin{proof}

	Seja

	[
	(x,y)\in A\times B.
	]

	Então

	[
	x\in A
	\quad\text{e}\quad
	y\in B.
	]

	Como

		[
			A\subseteq C
			\quad\text{e}\quad
			B\subseteq D,
		]

	segue que

		[
			x\in C
			\quad\text{e}\quad
			y\in D.
		]

	Portanto

		[
			(x,y)\in C\times D.
		]

\end{proof}

\begin{prop}[Distribuição sobre a União]
	\label{prop:matematica:teoria-conjuntos:produto-cartesiano:distribuicao-uniao}

	Para quaisquer conjuntos (A), (B) e (C),

	[
			A\times(B\cup C)
			=
			(A\times B)
			\cup
			(A\times C).
		]

\end{prop}

\begin{prop}[Distribuição sobre a Interseção]
	\label{prop:matematica:teoria-conjuntos:produto-cartesiano:distribuicao-intersecao}

	Para quaisquer conjuntos (A), (B) e (C),

	[
			A\times(B\cap C)
			=
			(A\times B)
			\cap
			(A\times C).
		]

\end{prop}

\begin{defn}[Produto Cartesiano Finito]
	\label{def:matematica:teoria-conjuntos:produto-cartesiano:produto-finito}

	Sejam

		[
			A_1,A_2,\dots,A_n
		]

	conjuntos.

	Define-se

	[
	A_1\times A_2\times\cdots\times A_n
	]

	como o conjunto de todas as (n)-uplas

	[
	(a_1,a_2,\dots,a_n)
	]

	tais que

		[
			a_k\in A_k
		]

	para todo

		[
			k=1,\dots,n.
		]

\end{defn}

\begin{defn}[Potência Cartesiana]
	\label{def:matematica:teoria-conjuntos:produto-cartesiano:potencia-cartesiana}

	Seja (A) um conjunto e (n\in\mathbb{N}).

	Define-se

	[
	A^n
	=
	\underbrace{
		A\times A\times\cdots\times A
	}_{n\text{ fatores}}.
	]

\end{defn}

\begin{ex}

	[
		\mathbb{R}^2
		=
		\mathbb{R}\times\mathbb{R}
	]

	é o plano cartesiano.

	Analogamente,

	[
			\mathbb{R}^3
		]

	representa o espaço tridimensional.

\end{ex}

\end{document}
